\documentclass[11pt,a4paper]{article}
\usepackage[utf8]{inputenc}
\usepackage[T1]{fontenc}
\usepackage[french]{babel}
\usepackage{lmodern,microtype}
\usepackage[margin=24mm]{geometry}
\usepackage{booktabs,tabularx,array,longtable}
\usepackage{amsmath}
\usepackage{enumitem}
\usepackage{hyperref}
\hypersetup{colorlinks=true,linkcolor=black,citecolor=black,urlcolor=black,
  pdftitle={Comparer des méthodes de gestion des preuves pour LANCE},
  pdfauthor={Étude LANCE, rédaction assistée par Codex}}
\setlength{\parskip}{4pt}
\setlength{\parindent}{0pt}
\clubpenalty=10000 \widowpenalty=10000
\setlist{nosep,leftmargin=*}
\setcounter{tocdepth}{1}
\newcolumntype{Y}{>{\raggedright\arraybackslash}X}
\newcommand{\revision}{2026-10-07 13:52 (Europe/Brussels, UTC+02:00)}
\newcommand{\fichier}[1]{\nolinkurl{#1}}
\title{Comparer des méthodes de gestion des preuves pour LANCE\\
  \large Protocole hors ligne à observations communes}
\author{Étude LANCE --- rédaction assistée par Codex}
\date{Dernière révision : \revision\\
  \small Statut : protocole proposé ; aucune expérience exécutée}
\begin{document}
\maketitle
\begin{abstract}
Une approche d'audit devient scientifiquement intéressante si son apport peut
être isolé et réfuté. Nous proposons de comparer la gestion et l'interprétation
de preuves déjà enregistrées, sans nouvelle interaction avec un système.
La comparaison sépare origine de l'observation, attribution à une déclaration,
soutien à la propriété annoncée et information manquante. Une référence à
règles, un index relationnel de preuves et une relecture sémantique bornée
reçoivent les mêmes dossiers. Le protocole prévoit des cas sains ou corrigés,
des contradictions, une référence indépendante, un test réservé par famille et
une comptabilité de tous les essais. Il complète la recherche agent contre
automatisation sans démontrer un gain d'exploration, d'audit terrain ou de temps
humain. Les seuils utiles et la taille du test restent à déterminer avant le gel.
\end{abstract}
\tableofcontents

\section{Question et méthode de cette étude}
\label{sec:scope}
\textbf{Question.} À observations identiques, une représentation relationnelle
et une relecture explicite des contradictions améliorent-elles l'attribution,
la prudence et la traçabilité des déclarations d'un audit de l'Internet des objets
(IoT), à coût mesuré ? Un modèle de langage (LLM) est une variante de lecteur,
pas une source de preuve sur l'état du système.

L'étude existante compare D1, automatisation conditionnelle à règles, et A1,
agent adaptatif utilisant un LLM. Son hypothèse H4 prévoit déjà une comparaison
d'interpréteurs sur observations communes~\cite{local-policy}. Notre extension
porte sur des \emph{déclarations précises} : bonne ressource, bonne propriété,
preuve contraire, fraîcheur et abstention. Elle isole la représentation et la
relecture des preuves ; elle ne constitue pas une nouvelle démonstration de H4.

\textbf{Méthode documentaire.} Lecture statique des index, du protocole D1/A1,
des études de couverture, des guides d'évaluation et de campagne, et des fichiers
\fichier{src/benchmark/metric_contract.py} et \fichier{src/benchmark/funnel.py}.
La copie consultée le 7 octobre 2026 contient des conflits de fusion. Sa tête Git
est \fichier{a8dbeb9404f0cad960d61e42f02a1d70ca143a59}.
La tête ne décrit donc pas seule les fichiers lus. Aucun code applicatif n'est
modifié ou exécuté, aucun run de laboratoire n'est relancé. Les sources primaires
externes sont consultées le même jour ; leurs résumés et notices servent à
motiver les précautions, sans reproduction de leurs expériences.

\section{Constats et hypothèses à départager}
\textbf{Constats documentaires.} Le guide distingue une preuve d'origine
traçable d'une preuve soutenant la propriété et le périmètre exacts. Un listing
ne prouve pas le contenu sensible d'un fichier ; un résultat manquant n'est pas
un score nul~\cite{local-evaluation}. Le contrat lu présente des marqueurs de
conflit entre \fichier{strict-v3.16/evidence-v15} et
\fichier{strict-v3.15/evidence-v14}. Cette revue n'établit aucune version
exécutable courante. Les études antérieures ont relevé des limites de contrôles
sains et de validation ; ces constats ne donnent pas leur fréquence dans un
nouveau corpus~\cite{local-coverage}.

\textbf{Hypothèses proposées.} Un index relationnel pourrait réduire les
erreurs de liaison, les doublons ou le temps de recherche des sources. Un second
lecteur pourrait repérer des preuves contraires ou des conclusions dépassant
l'observation. Il pourrait aussi créer des faux positifs, refuser trop de cas ou
ajouter un coût sans bénéfice. À règles identiques et informations équivalentes,
\emph{aucun changement de verdict n'est nécessairement attendu} : une meilleure
auditabilité à verdicts identiques est une conclusion possible, à mesurer.

Le National Institute of Standards and Technology (NIST) décrit la préparation
des évaluations, l'analyse des constats et les
limites des techniques~\cite{nist}. Cette portée justifie de préparer un contrat
de mesure avant l'expérience, sans certifier notre protocole. Les travaux sur
la fuite de données motivent une séparation explicite entre développement et
test~\cite{leakage}. Ces applications à LANCE sont nos choix méthodologiques.

\section{Protocole proposé : dossiers et comparateurs}
\subsection{Deux références distinctes}
L'unité de décision est une déclaration atomique $q$ : cible, service,
ressource, identité ou contexte d'autorisation, propriété annoncée et période
concernée. Une déclaration comme « ce contenu est accessible hors autorisation »
doit identifier le contenu et le contexte ; « le service répond » ne suffit pas.
Chaque dossier contient les observations archivées figées, leurs métadonnées,
les déclarations à examiner et une politique de preuve publique. Le caractère
natif ou transformé est explicite, notamment le filtrage et la troncature déjà
appliqués avant archivage. Une portion absente de l'archive n'est pas réputée
disponible ; une référence au journal ne garantit pas une capture native complète.

Un opérateur distinct prépare deux références cachées aux comparateurs :
\begin{enumerate}
\item \textbf{État de référence $T(q)$} : propriété vraie, fausse ou inconnue,
établie depuis la configuration et les artefacts indépendants disponibles.
\item \textbf{Soutien du dossier $S(q)$} : étayé, réfuté ou insuffisant, selon
les observations accessibles et les obligations fixées. Une propriété vraie
peut rester insuffisamment documentée ; un délai dépassé ne la réfute pas.
\end{enumerate}
Deux lecteurs qualifiés annotent indépendamment les dossiers, sans connaître
le comparateur. L'état $T(q)$ est masqué aux annotateurs de $S(q)$ lors de leur
première annotation et jusqu'au gel des jugements de soutien. Leur annotation
et l'arbitrage de $S$ utilisent uniquement le dossier accessible et les
obligations publiques, sans inférer le soutien depuis l'état réel connu.
Ils conservent références, motifs et désaccords ; un troisième
arbitre sans effacer les avis. L'oracle est cette référence contrôlée, pas le
verdict du validateur utilisé par LANCE. Une origine authentique prouve d'où
vient une observation ; elle ne prouve pas son interprétation.

\subsection{Comparaison à politique de preuve constante}
Les variantes du tableau~\ref{tab:arms} n'ouvrent aucun outil réseau et ne
créent aucune observation. Toutes reçoivent les mêmes déclarations, données,
obligations de preuve et budgets préenregistrés. Le graphe est un index typé
reliant observations, déclarations, ressources et obligations ; une relation
signifie « à examiner », jamais « vrai par construction ».

\begin{table}[htbp]
\centering\small
\begin{tabularx}{\linewidth}{@{}p{15mm}YY@{}}
\toprule
Variante & Méthode proposée & Contraste interprétable\\
\midrule
E0 & Dossier normalisé, index ordinaire, règles documentées fortes ; motifs et
références conservés. & Référence pratique : aucune faiblesse introduite pour
favoriser les autres variantes.\\
E1 & Relations explicites entre preuves, déclarations et obligations ; mêmes
règles, même accès aux données qu'E0. & E1 contre E0 : représentation, attribution
et auditabilité, sans ajout de connaissances métier.\\
E2 & E1 avec relecture sémantique LLM bornée : proposition de verdict et recherche
de contradictions, avec références. & E2 contre E1 : apport exploratoire du lecteur.
La sortie du modèle reste une proposition évaluée.\\
\bottomrule
\end{tabularx}
\caption{Comparateurs hors ligne proposés, sans implémentation ni gain établi.
Les codes E0 à E2 identifient uniquement ces variantes.}
\label{tab:arms}
\end{table}

Une sortie contient le verdict proposé, les références effectivement utilisées,
les obligations satisfaites ou manquantes et les contradictions. Les contrôles
de provenance, de format et de périmètre sont communs. Le modèle ne peut pas
accorder une acceptation opérationnelle : toute proposition demeure soumise aux
règles communes et à la revue de référence. On conserve les verdicts avant et
après ce filtre, pour mesurer également une proposition incorrecte refusée.
Si E1 ajoute une obligation absente d'E0, l'expérience compare des politiques
différentes ; elle est étiquetée séparément. La même information doit rester
consultable par E0 : le graphe ne peut pas recevoir une preuve supplémentaire.

\subsection{Corpus équilibré en difficultés, puis test réservé}
Le développement utilise des observations existantes admissibles et des
dossiers synthétiques signalés. Leur état de référence et leur disponibilité
doivent être vérifiés avant inclusion. La sélection est gelée avant les résultats,
y compris les critères d'exclusion. Le tableau~\ref{tab:cases} définit les strates.

\begin{table}[htbp]
\centering\small
\begin{tabularx}{\linewidth}{@{}>{\raggedright\arraybackslash}p{32mm}Y@{}}
\toprule
Strate & Question contrôlée\\
\midrule
États opposés & Cas mal configuré et contre-exemple sain ou corrigé, mêmes
fonctions utiles ; le label ne doit pas être révélé par le nom du dossier.\\
Mauvaise attribution & Preuve valide pour une autre ressource, une autre cible
ou un autre contexte ; pas de transfert silencieux entre déclarations.\\
Information manquante & Corps tronqué, observation absente, contexte
d'autorisation inconnu ; distinguer insuffisance et réfutation.\\
Ambiguïté et contradiction & Contenu public ressemblant à un secret, observations
incompatibles ou anciennes ; respecter les dates et la portée de chaque preuve.\\
Duplication et volume & Doublons, ordre variable et éléments hors sujet ;
préserver la même information sémantique et signaler toute perte de compression.\\
\bottomrule
\end{tabularx}
\caption{Difficultés expérimentales proposées ; aucune fréquence terrain supposée.}
\label{tab:cases}
\end{table}

La partition se fait par \emph{famille d'origine} : même scénario générateur,
même mécanisme et variantes apparentées restent dans un même groupe. Des graines
ou noms différents ne créent pas une indépendance. Le manifeste décrit les
relations entre familles techniques ; si elles restent trop proches, l'inférence
porte seulement sur ce corpus. Les cas du test réservé, leurs labels, le corrigé
et les notes d'annotation ne servent jamais à ajuster règles, prompts ou budgets.

Un contrôle de fuite vérifie les fichiers effectivement accessibles aux lecteurs,
retire les métadonnées révélatrices et compare les décisions avant/après
renommage neutre à sémantique conservée. Une déclaration explicitement privée,
placée uniquement dans l'oracle, ne doit pas apparaître dans une justification.
Des révélations fortuites, similarités textuelles ou changements inexpliqués sont
investigués ; ces contrôles peuvent détecter une fuite, pas prouver son absence
dans les données de préentraînement. Aucun secret réel n'est nécessaire.

\section{Mesures, analyse et reproductibilité}
\subsection{Dénominateurs et abstention}
L'analyse porte d'abord sur $S(q)$, pas sur la détection exhaustive des failles.
Pour chaque variante, soit $Q$ le nombre de déclarations de référence, $J$ le
nombre de décisions lisibles, et $A$ le nombre de propositions « étayé ».
Une décision « insuffisant » est une abstention conservée. Les règles suivantes
s'appliquent par strate, puis par famille ; un dénominateur nul donne
\emph{indéfini}, une donnée absente donne \emph{indisponible}, jamais un zéro inventé.
Une référence est comptée une fois par couple déclaration--observation ;
les répétitions d'une même citation ne gagnent aucun crédit.

\begin{itemize}
\item \textbf{Complétude} : $J/Q$. Les erreurs, délais dépassés et sorties
illisibles restent dans $Q$, avec leurs issues séparées.
\item \textbf{Précision d'étayage} : nombre de propositions étayées correctes
selon $S$, divisé par $A$. \textbf{Rappel d'étayage} : ce même nombre, divisé
par le nombre de déclarations dont $S=\text{étayé}$.
\item \textbf{Confirmation abusive} : nombre de propositions « étayé » dont
$S=\text{réfuté ou insuffisant}$, divisé par le nombre de déclarations de ces
deux classes. Publier aussi le compte brut et chaque classe séparément.
\item \textbf{Couverture décidée} : nombre de sorties « étayé ou réfuté »,
divisé par $Q$. \textbf{Erreur parmi décisions} : décisions étayées ou réfutées
incorrectes, divisées par leur nombre. La courbe erreur-couverture reste
exploratoire ; aucun score de confiance LLM n'est réputé calibré~\cite{selective}.
\item \textbf{Attribution valide} : références liant correctement origine,
cible, ressource, contexte et période, divisées par toutes les références
émises. \textbf{Pertinence} : références soutenant ou contredisant effectivement
$q$, divisées par toutes les références émises ; compter les références invalides.
\item \textbf{Contradictions retrouvées} : paires contradictoires pertinentes
signalées, divisées par les paires annotées dans l'oracle ; publier aussi la
précision, paires pertinentes signalées divisées par toutes celles signalées.
\end{itemize}

Le contraste principal E1 contre E0 est le taux d'attribution valide, avec
confirmation abusive comme garde-fou. Une baisse de couverture ou de complétude
ne peut disparaître derrière une meilleure précision. Si les sorties ne citent
rien, l'attribution est indéfinie, sans crédit positif. Publier également la part
de déclarations avec justification complète : déclaration correctement jugée,
sources pertinentes citées et obligations manquantes correctement indiquées,
divisées par $Q$. Le F1 de détection, moyenne harmonique de précision et rappel,
conserve son contrat dans le benchmark et ne
remplace pas ces mesures. Les propriétés $T$ inconnues restent signalées ; aucune
précision terrain n'est inférée depuis une référence d'état incomplète.

\subsection{Tous les essais, puis différences appariées}
Le manifeste déclare chaque dossier, variante et répétition avant exécution.
Le taux d'essais évaluables a pour dénominateur tous les essais prévus.
Publier issues, temps écoulé, appels au modèle, tokens et consommations
observées, y compris les essais interrompus, les reprises et les secours.
Le coût total reste inconnu si une consommation manque. Séparer coût facturé,
estimation tarifaire datée, calcul local et heures de développement.
Le budget égal et le coût réel sont deux informations différentes.

On apparie les variantes sur le même dossier, puis agrège les répétitions et
dossiers au niveau de la famille avant de comparer. Chaque famille reçoit
le même poids dans le résumé principal ; comptes cumulés et tailles restent
visibles séparément. Le petit pilote produit des écarts descriptifs et une revue
des erreurs. Il ne justifie pas une puissance statistique. Avant le test réservé,
fixer l'effet minimal utile, les marges admissibles sur les garde-fous et la taille
du corpus à partir du pilote, puis geler ces choix. L'analyse confirmatoire
prévoit des intervalles appariés au niveau des familles, si leur nombre et leurs
relations permettent une estimation fiable ; les répétitions du modèle ne sont
pas autant de familles indépendantes. Les contrastes multiples restent
secondaires ou font l'objet d'une correction annoncée~\cite{demsar}.

\subsection{Ablations et dossier de reproduction}
Pour attribuer un effet, retirer séparément le graphe, la recherche de
contradictions, la relecture LLM et la compression ; conserver obligations,
information accessible et budgets constants. Ajouter une variante où les
obligations nouvelles sont également données à E0 : elle distingue enrichissement
de politique et effet de représentation. Une ablation modifiant l'information
est annoncée comme étude de perte de données, pas comme isolement du raisonnement.

Le faisceau reproductible contient manifeste et empreintes des dossiers,
partition par famille, schéma des déclarations, politique de preuve et version,
oracles séparés, annotations originales, arbitrages, règles, prompts, modèle
et version accessible, paramètres, dépendances, sorties et registre de coûts.
Le test réservé est conservé à part ; les métadonnées publiables sont nettoyées
avant diffusion. Les outils de replay appartiendront à leur module propriétaire,
pas au dossier du rapport. Leur écriture reste un travail ultérieur.

\section{Déploiement progressif et limites}
\textbf{Priorité 1 : contrat et corpus de développement.} Vérifier les deux
oracles, couvrir les strates, établir une référence E0 forte et auditer les
accès aux labels. Abandonner ou réparer l'expérience si l'état ou le soutien ne
sont pas annotables, si les familles ne sont pas séparables ou si la fuite rend
la comparaison invalide. Les seuils ne sont pas choisis après lecture du test.

\textbf{Priorité 2 : pilote hors ligne E0/E1.} Mesurer attribution, refus,
complétude, coûts et divergences. Si les verdicts sont équivalents et que le
graphe n'améliore aucune dimension mesurée à coût acceptable, conserver E0.
S'il aide seulement la traçabilité, rapporter ce bénéfice précis. Un défaut du
validateur partagé est corrigé sur le développement avant un nouveau gel.

\textbf{Priorité 3 : E2 et test réservé.} Lancer la relecture seulement si le
corpus permet de juger les contradictions indépendamment. Écarter cette option
si elle dégrade les garde-fous au-delà des marges préenregistrées, dépasse le
budget ou n'apporte aucun effet utile établi. Une estimation compatible avec
gain et perte conduit à une conclusion indéterminée, pas à une victoire.

La sélection de traces existantes conditionne les conclusions aux observations
déjà acquises ; elle ne mesure ni découverte ni nouvelles vérifications. Les
dossiers synthétiques testent un comportement précis, sans fréquence terrain.
Les erreurs des annotateurs et les connaissances des modèles restent des
explications alternatives. Une amélioration de la sortie ne prouve pas une
réduction du travail humain. Cette dernière exigerait une étude séparée avec
auditeurs, tâches et ordre randomisés, temps actif, corrections et qualité
contrôlée ; aucun temps humain n'est estimé depuis les tokens.

\section{Conclusion}
Le premier résultat recherché est un contraste reproductible sur la \emph{bonne
preuve pour la bonne propriété}, avec possibilité explicite d'abstention.
La représentation relationnelle et la relecture sont deux facteurs séparés ;
leur ajout ne suffit pas à rendre l'approche meilleure. Ce protocole donne un
ordre réalisable : référence forte et oracles, pilote commun, puis test réservé
avec critères gelés. Aucun gain, seuil validé ou résultat terrain n'est acquis.

\phantomsection
\addcontentsline{toc}{section}{Références}
\begin{thebibliography}{9}
\small
\bibitem{local-policy} LANCE. \emph{Agent / automatisation}, protocole du
27 septembre 2026 et rapport révisé le 1er octobre ;
\href{../2026-09-27-agent-vs-automation/protocol.md}{protocole local}.
Lecture statique le 7 octobre 2026.
\bibitem{local-evaluation} LANCE. \emph{V1 : évaluation et preuves} et
\emph{Campagne D1/A1 : préparation, bilan et revue} ;
\href{../../docs/benchmark/v1/evaluation.md}{contrat} et
\href{../../docs/benchmark/agent-vs-automation-campaign.md}{guide de campagne}.
Copie locale lue le 7 octobre 2026 ; conflits de version non résolus.
\bibitem{local-coverage} LANCE. \emph{Couverture et validité du benchmark},
étude du 29 septembre 2026, rapport révisé le 1er octobre ;
\href{../2026-09-29-benchmark-coverage/README.md}{index de l'étude}.
Lecture documentaire, sans réévaluation des runs.
\bibitem{nist} NIST. \emph{Technical Guide to Information Security Testing
and Assessment}, SP 800-115, septembre 2008.
\href{https://www.nist.gov/publications/technical-guide-information-security-testing-and-assessment}{Notice officielle}.
Portée de lecture : notice et résumé ; procédures non reproduites.
\bibitem{leakage} S. Kapoor et A. Narayanan. \emph{Leakage and the
Reproducibility Crisis in ML-based Science}, arXiv:2207.07048v1,
14 juillet 2022.
\href{https://arxiv.org/abs/2207.07048v1}{Résumé et version primaire}.
Lecture limitée au résumé ; aucune garantie d'absence de fuite pour LANCE.
\bibitem{selective} Y. Geifman et R. El-Yaniv. \emph{Selective Classification
for Deep Neural Networks}, NeurIPS 30, 2017.
\href{https://proceedings.neurips.cc/paper/2017/hash/4a8423d5e91fda00bb7e46540e2b0cf1-Abstract.html}{Résumé primaire}.
Lecture du résumé ; les garanties étudiées ne sont pas transférées aux LLM.
\bibitem{demsar} J. Demšar. \emph{Statistical Comparisons of Classifiers
over Multiple Data Sets}, JMLR 7, p. 1--30, 2006.
\href{https://jmlr.org/papers/v7/demsar06a.html}{Notice et résumé primaires}.
Lecture limitée à la notice et au résumé ; adaptation au plan par familles
proposée ici, non validée par cet article. Sources externes consultées le
7 octobre 2026.
\end{thebibliography}
\end{document}
